\section*{Chapter \ref{ch:g9:demand-and-supply-equations} --- Demand and Supply with Equations}

\begin{solution}{exo:g9:demand-and-supply-equations:1}
2000: real \omterm{def:g7:money:price}{price} 100, $48.8$ subscriptions per 100 people. 2010: real \omterm{def:g7:money:price}{price} $9.3$, $91.1$ subscriptions
per 100 people.
\end{solution}

\begin{solution}{exo:g9:demand-and-supply-equations:2}
$200 - 5 \times 12 = 140$ units.
\end{solution}

\begin{solution}{exo:g9:demand-and-supply-equations:3}
$90 - 3p = 10 + 5p$, so $80 = 8p$ and $p^* = 10$ dalens; $Q^* = 90 - 30 = 60$ units.
\end{solution}

\begin{solution}{exo:g9:demand-and-supply-equations:4}
The intercept is 300; the quantity falls by 100 for 10 dalens, that is 10 per dalen: $Q^d = 300 - 10p$.
\end{solution}

\begin{solution}{exo:g9:demand-and-supply-equations:5}
$150 - 4p = 20 + 6p$, so $130 = 10p$ and $p^* = 13$ dalens; $Q^* = 150 - 52 = 98$ units.
\end{solution}

\begin{solution}{exo:g9:demand-and-supply-equations:6}
$120 - 4p = 6p$, so $p^* = 12$ dalens; $Q^* = 120 - 48 = 72$ units.
\end{solution}

\begin{solution}{exo:g9:demand-and-supply-equations:7}
$Q^s = 32 + 6p$. Then $120 - 4p = 32 + 6p$, so $88 = 10p$ and $p^* = 8.8$ dalens (quantity $120 - 35.2 =
84.8$ units). The \omterm{def:g7:money:price}{price} fell by $10 - 8.8 = 1.2$ dalens, not by the full 2 dalens: sellers keep the
rest.
\end{solution}

\begin{solution}{exo:g9:demand-and-supply-equations:8}
Demanded: $120 - 32 = 88$ units; supplied: $20 + 48 = 68$ units; \omterm{def:g7:markets-and-prices:shortage}{shortage}: $88 - 68 = 20$ units.
\end{solution}

\begin{solution}{exo:g9:demand-and-supply-equations:9}
A rise in supply alone lowers the \omterm{def:g7:money:price}{price}, so a higher \omterm{def:g7:money:price}{price} does not fit a supply rise alone. Something
else must have changed: for example, demand rose at the same time, by more than enough to lift the
\omterm{def:g7:money:price}{price}.
\end{solution}

\begin{solution}{exo:g9:demand-and-supply-equations:10}
$150 - 4p = 40 + 6p$, so $110 = 10p$ and $p^* = 11$ dalens; $Q^* = 150 - 44 = 106$ units. Both \omterm{def:g9:demand-and-supply-equations:shock}{shocks}
raise the quantity, so its rise could be predicted without numbers; they push the \omterm{def:g7:money:price}{price} in opposite
directions, so the \omterm{def:g7:money:price}{price} change needs the numbers (here it rose from 10 to 11).
\end{solution}

\begin{solution}{exo:g9:demand-and-supply-equations:11}
New supply: $Q^s = 32 + 6p$. With $Q^d = 50 - p$: before, $30 = 7p$, $p^* \approx 4.29$; after, $18 = 7p$,
$p^* \approx 2.57$: a fall of $12 \div 7 \approx 1.71$ dalens. With $Q^d = 200 - 12p$: before, $180 =
18p$, $p^* = 10$; after, $168 = 18p$, $p^* \approx 9.33$: a fall of $12 \div 18 \approx 0.67$ dalens.
Adding 12 units at every \omterm{def:g7:money:price}{price} is the same as a cost fall of $12 \div 6 = 2$ dalens per unit. Buyers
get most of it in the first \omterm{def:g7:markets-and-prices:market}{market}, where they hardly change their purchases when the \omterm{def:g7:money:price}{price} changes
(steep demand).
\end{solution}

\begin{solution}{exo:g9:demand-and-supply-equations:12}
Each point of the chart is one equilibrium, a crossing of the two curves that year. The chart cannot
say whether supply moved, demand moved, or both. Subscriptions stopped rising while the \omterm{def:g7:money:price}{price} kept
falling: demand may have stopped growing because nearly everyone already had a phone, but the chart
alone cannot show it.
\end{solution}

\begin{solution}{pb:g9:demand-and-supply-equations:1}
\textbf{1.} 100 kg less for each 2 dalens: 50 kg per dalen.
\textbf{2.} At $p = 2$, $Q^d = 900 = a - 50 \times 2$, so $a = 1\,000$: $Q^d = 1\,000 - 50p$.
\textbf{3.} The \omterm{def:g7:demand-and-supply-first-look:quantity-supplied}{quantity supplied} rises by 150 kg for each 2 dalens: $d = 75$; at $p = 2$, $200 = c +
150$, so $c = 50$: $Q^s = 50 + 75p$.
\textbf{4.} Demanded: $1\,000 - 250 = 750$ kg; supplied: $50 + 375 = 425$ kg.
\textbf{5.} $1\,000 - 50p = 50 + 75p$, so $950 = 125p$ and $p^* = 7.6$ dalens per kg; $Q^* = 1\,000 - 380
= 620$ kg.
\textbf{6.} $50 + 75 \times 7.6 = 50 + 570 = 620$ kg.
\textbf{7.} Demanded 500 kg, supplied 800 kg: a \omterm{def:g7:markets-and-prices:market-surplus}{market surplus} of 300 kg.
\textbf{8.} Demanded 700 kg, supplied 500 kg: a \omterm{def:g7:markets-and-prices:shortage}{shortage} of 200 kg.
\textbf{9.} $7.6 \times 620 = 4\,712$ dalens.
\textbf{10.} $175 + 75 \times 6 = 625$ kg.
\textbf{11.} $1\,000 - 50p = 175 + 75p$, so $825 = 125p$ and $p^* = 6.6$ dalens per kg; $Q^* = 1\,000 -
330 = 670$ kg.
\textbf{12.} The \omterm{def:g7:money:price}{price} fell by 1 dalen per kg (from $7.6$ to $6.6$); the quantity rose by 50 kg (from 620
to 670).
\textbf{13.} $6.6 \times 670 = 4\,422$ dalens: it fell by $4\,712 - 4\,422 = 290$ dalens, although they
sell more fish.
\end{solution}
